% Generated draft: requires mathematical and source review.
\documentclass[11pt]{article}
\usepackage[margin=25mm]{geometry}
\usepackage{fontspec}
\setmainfont{Latin Modern Roman}
\usepackage{amsmath,amssymb,mathtools,unicode-math}
\setmathfont{Latin Modern Math}
\usepackage{enumitem,needspace,xurl,hyperref}
\setlength{\parindent}{0pt}
\setlength{\parskip}{5pt}
\setlength{\emergencystretch}{3em}
\newcommand{\N}{\mathbb{N}}
\newcommand{\Z}{\mathbb{Z}}
\newcommand{\Q}{\mathbb{Q}}
\newcommand{\R}{\mathbb{R}}
\newcommand{\C}{\mathbb{C}}
\newcommand{\F}{\mathbb{F}}
\newcommand{\A}{\mathbb{A}}
\newcommand{\PP}{\mathbb{P}}
\newcommand{\E}{\mathbb{E}}
\newcommand{\eps}{\varepsilon}
\newcommand{\dd}{\,\mathrm{d}}
\begin{document}
\section*{UnsolvedMath ID: 801008}
\textbf{Source problem code:} catalogue entry
\subsection{Definitions and mathematical statement}
Let \(\beta\) be an ordinal and let \(L_\beta\) denote the corresponding level of the constructible hierarchy. For a set-theoretic structure \(M\), \(\mathrm{HC}^{M}\) denotes its internally hereditarily countable sets: those sets whose transitive closures are at most countable as computed in \(M\). Thus \(L_\beta\models V=\mathrm{HC}\) means that every set in \(L_\beta\) is hereditarily countable according to \(L_\beta\).

Write \(\mathcal R_\beta=\mathcal P(\omega)\cap L_\beta\). Under the usual second-order-arithmetic interpretation, satisfaction by \(\mathcal R_\beta\) refers to the structure with number domain \(\omega\), set domain \(\mathcal R_\beta\), and the usual arithmetic operations and membership relation. The supplied record does not explicitly specify this interpretation or the background assumptions ensuring it applies.

The source asks whether the following equivalence holds, subject to its unspecified conventions for the index \(m\):
\[
L_\beta\models V=\mathrm{HC}
\quad\Longrightarrow\quad
\left[
\mathcal R_\beta\models\Pi^1_{m+1}\mathsf{-MI}
\quad\Longleftrightarrow\quad
\beta\text{ supports a strong }\Sigma_{m+1}\text{-nesting}
\right].
\]
Here \(\Pi^1_{m+1}\mathsf{-MI}\) is a provisional rendering of the source's malformed expression \texttt{\textbackslash mathbf\{s\}ymbol\{\textbackslash Pi\}} followed by the superscript \(1\), subscript \(m+1\), and \(\mathsf{-MI}\). The abbreviation presumably refers to a monotone-induction principle, but its exact schema, permitted parameters, and lightface or boldface convention are not supplied. The source also provides no definition of a strong \(\Sigma_{m+1}\)-nesting or of what it means for \(\beta\) to support one. Consequently, the displayed equivalence records the proposed target but is not a fully specified mathematical assertion until these conventions, the range of \(m\), and any additional hypotheses on \(\beta\) are recovered.
\subsection{Short English statement}
The supplied problem asks whether, assuming \(L_\beta\models V=\mathrm{HC}\), the reals in \(L_\beta\) satisfy the indicated projective monotone-induction principle exactly when \(\beta\) supports a strong \(\Sigma_{m+1}\)-nesting. The induction principle, nesting condition, and index range require clarification. The registry labels the problem open; that label is supplied metadata, not an independently verified assessment.
\subsection{Sources}
\begin{itemize}
\item Local UnsolvedMath registry, permanent ID \(801008\), exact source problem code \texttt{catalogue entry}, title \emph{Open Problems in Mathematical Logic --- Problem 20}; collection: AMR Open Problem Lists. The registry attributes the record to Barmpalias, Gao et al., with proposed year 2026 and supplied status \texttt{open}.
\item Supplied source-list reference: Barmpalias, Gao et al., \emph{Open Problems in Mathematical Logic} (2026), Problem 20, \url{https://arxiv.org/abs/2608.26628}. The background records access on 2026-10-16 and extraction by \texttt{arxiv\_latex}; these are supplied provenance claims, not independent source verification.
\item Supplied upstream identity: \(801008\), revision \texttt{d371ee917ac9647fbb06cd6fc529deec6f3dda82}. The registry record was created and updated at \texttt{2026-10-16T14:00:00Z}.
\end{itemize}
\end{document}
